\section{The Household Population}
\label{ch:population}

We construct a synthetic population of 5{,}000 household agents from South African microdata
This section outlines the data sources, variable derivations, financial indicator imputations, and empirical validation of the static population structure.
Dynamic household interactions and behavioural realism are evaluated separately in Section~\ref{sec:results-baseline}.
All monetary figures are denominated in 2017 South African Rand (ZAR) matching the base survey year, with no inflation adjustments applied at the data layer.

\subsection{Primary Dataset: NIDS Wave 5}
\label{sec:backbone}

The backbone of the synthetic population is drawn from \ac{NIDS} Wave 5 \cite{nids_w5_2018}.
Filtering the raw sample ($N = 13,719$) for households with reported income, valid size, and
positive sampling weights leaves 10{,}841 records. The 2{,}878 discarded records carry zero or missing weights; their removal thus entails no loss of weighted population mass. Scaled by national design weights, the final sample represents 18.67 million households or roughly 56.5 million individuals.

We derive six core economic variables per household:
\begin{itemize}
    \item \textbf{Monthly income:} Total aggregated household income as derived by \ac{NIDS}.
    \item \textbf{Income source:} Primary income stream, categorised into wage, government grant, or remittance/other.
    \item \textbf{Committed expenditure:} Sum of food outlays and paid rent, establishing an incompressible consumption floor.
    \item \textbf{Discretionary expenditure:} Non-food expenditure (recorded in \ac{NIDS} independently of rent).
    \item \textbf{Liquid savings:} Proxied by self-reported financial assets and winsorised at the 99th percentile (ZAR 133{,}662) to limit tail-distorted variance.
    \item \textbf{Traditional debt:} Total outstanding financial debt reported in \ac{NIDS}.
\end{itemize}

Total expenditure in \ac{NIDS} decomposes into food, non-food, rent paid, and imputed rent.
We exclude imputed rent, the non-cash consumption value of ower-occupied housing, from expenditure tiers.
Including it would artificially inflate compressible consumption by a median of ZAR 700 per month across the 82.6\% of households that own their homes.

Four survey fields record values only when an asset or liability is present, omitting explicit zeros.
Missing entries are treated as non-holding (zero balances): financial assets (unobserved in 48.2\% of cases), financial debt (56.6\%), total debt (54.4\%), and rent paid (80.2\%).
We test this non-response assumption against aggregate national credit data in Section~\ref{sec:data-validation}.
Operationally, this assumption implies that committed expenditure reduces to food costs alone for homeowning households.

Households are stratified into per-capita income quintiles bounded at ZAR 900, ZAR 1{,}801, ZAR 3{,}400, and ZAR 7{,}712 per month.
Demographic variables for the household head (age, gender, race, education) are merged from the \ac{NIDS} individual file, matching 10{,}736 of the 10{,}841 retained households.

\subsection{Imputing Financial Inclusion: FinScope 2019}
\label{sec:match}

Because \ac{NIDS} omits institutional banking status and specific debt instruments, we impute six binary financial inclusion indicators from the 2019 \ac{FinScope} survey \cite{finscope2019}.
Banked status determines entry into \ac{BNPL}, bounding potential market adoption at the observed 82.8\% macro limit rather than an arbitrary baseline parameter.
Five credit-product flags (store cards, revolving credit, hire purchase, short-term loans, and personal loans) determine the interest rates and repayment schedules assigned in Section~\ref{sec:servicing}.
Unused indicator fields are omitted.

We transfer attributes using cell-donor matching.
Stratifying both datasets by per-capita income quintile and province, each \ac{NIDS} record draws a donor from its matching cell with probability proportional to the donor's survey weight (seed: 42).
Extracting all six indicators from a single donor preserves the joint product distribution.
Each of the 45 sampling cells contains at least 30 donor candidates.

\ac{FinScope} reports income in seven discrete bands rather than continuously; we place donors into quintiles using band midpoints divided by household size.
Deflating 2019 income brackets to 2017 Rands shifts imputed quintile rates by at most 2.4 percentage points, so we retain the unadjusted key.

Table~\ref{tab:imputation_errors} shows national target matches, where the maximum absolute discrepancy is 1.8 percentage points (store cards).
Figure~\ref{fig:inclusion-fidelity} breaks down alignment across quintiles; the maximum deviation across all thirty quintile-by-indicator cells is 2.3 percentage points.

\begin{table}[H]
\centering
\caption{National rates for the imputed indicators: the FinScope 2019 benchmark, the weighted synthetic population carrying the flags, and the unweighted 5{,}000 agent population. The difference column reports the synthetic population less FinScope, and is computed before rounding.}
\label{tab:imputation_errors}
\begin{tabular}{lrrrr}
\toprule
\textbf{Imputed indicator} & \textbf{FinScope 2019} & \textbf{Synthetic} & \textbf{Agents} & \textbf{Difference} \\
 & \textbf{(\%)} & \textbf{(\%)} & \textbf{(\%)} & \textbf{(pp)} \\
\midrule
Banked status    & 82.4 & 82.8 & 82.8 & 0.4 \\
Store card       & 20.8 & 22.6 & 21.8 & 1.8 \\
Revolving credit &  3.7 &  4.1 &  3.8 & 0.5 \\
Hire purchase    &  6.5 &  6.9 &  6.5 & 0.5 \\
Short-term loan  &  2.7 &  3.4 &  3.3 & 0.6 \\
Personal loan    &  2.1 &  2.3 &  2.3 & 0.2 \\
\bottomrule
\end{tabular}
\end{table}

\begin{figure}[htbp]
\centering
\includegraphics[width=\textwidth]{figures/fig_inclusion_fidelity.pdf}
\caption{Fidelity of the imputed indicators. Panel (a) compares national rates across the FinScope 2019 benchmark, the weighted synthetic population, and the unweighted agent population. Panel (b) decomposes the synthetic-to-FinScope gap by income quintile; the largest single deviation is 2.3 percentage points, against a tolerance of 5 percentage points.}
\label{fig:inclusion-fidelity}
\end{figure}

The geographic and income-based cells establish the reporting categories for quintile-level analysis and define the reference networks that drive peer-influence dynamics in the simulation.

\subsection{Debt Servicing Construction}
\label{sec:servicing}

\ac{NIDS} measures debt \emph{stocks}, but the simulation requires monthly cash \emph{flows}.
We reconstruct monthly debt service through a three-step amortization procedure using statutory maximum rates and empirical repayment horizons.

First, outstanding debt balances are amortized using the statutory ceiling rates mandated by the \ac{NCA} in 2016 \cite{sa_ncr_ratecaps_2016}.
Relative to the 2017 repo rate (7.00\%), these caps sit at 21\% for credit facilities, 24\% for standard credit agreements, and 28\% for unsecured credit.
Because realized average \acp{APR} by product class are unobserved in \ac{CCMR} reports \cite{ncr_ccmr_2017} or collateral 2017 sources, statutory caps serve as proxies.
This sets an upper bound on interest rates, introducing a systematic upward bias to calculated debt service flows.

Second, product classes are assigned specific repayment horizons (Table~\ref{tab:credit-terms}). Horizons for store cards and hire purchase equal the median ratio of total debt balance to observed monthly payment in \ac{NIDS}.
Because the numerator aggregates overall financial debt while the denominator isolates single product payments, these derived horizons are conservative upper bounds.
For other categories, horizons match the \ac{CCMR} stock-flow implied duration (calculated as three times the gross debtors book divided by quarterly credit granted).
This proxy tracks official \ac{CCMR} agreement term distributions within roughly three months.

Because \ac{NIDS} aggregates debt into a single household-level stock, rate and horizon parameters cannot be weighted by individual loan balances.
Instead, households adopt the unweighted mean interest rate and horizon across the product classes flagged by their donor.
Households whose donor holds no flagged product are priced under unsecured terms.

\begin{table}[H]
\centering
\caption{Repayment horizons by product class, with provenance. Interest rates are the \ac{NCA} statutory maxima at the 2017 repurchase rate throughout.}
\label{tab:credit-terms}
\begin{tabular}{L{2.7cm}L{1.2cm}L{1.9cm}L{5.7cm}}
\toprule
\textbf{Product class} & \textbf{\ac{APR}} & \textbf{Horizon (months)} & \textbf{Provenance} \\
\midrule
Store card & 21\% & 11 & Measured in \ac{NIDS}: median balance over observed monthly clothing-account payment, $n=782$ \cite{nids_w5_2018} \\
\addlinespace
Hire purchase & 24\% & 8 & Measured in \ac{NIDS}: median balance over observed monthly hire-purchase payment, $n=121$ \cite{nids_w5_2018} \\
\addlinespace
Short-term loan & 60\% & 1 & \ac{FinScope} defines the product as repayable within 31 days; the \ac{CCMR} places 65.8\% of such agreements at one month or less \cite{finscope2019, ncr_ccmr_2017} \\
\addlinespace
Personal loan & 28\% & 25 & \ac{CCMR} stock-flow implied life of the unsecured book, 24.8 months \cite{ncr_ccmr_2017} \\
\addlinespace
Revolving credit & 21\% & 20 & No contractual term exists. Derived as the reciprocal of the mid-range minimum-payment fraction, which is swept over 0.025 to 0.10 \\
\addlinespace
Unflagged balance & 28\% & 25 & Fallback for the 74.0\% of donors flagging no specific product; mapped as personal loan \\
\bottomrule
\end{tabular}
\end{table}

Third, generated monthly repayments are capped using the affordability calculation under \ac{NCA} Regulation 23A.
The regulation mandates minimum living expense allowances by gross income band, treating only the remaining balance as disposable for debt service.
Applied to household-level income, the maximum allowable debt service ceiling ranges from zero to 93.2\% of gross income.
Applying an individual legal framework to household income relaxes borrowing constraints for large, low-income households in the bottom quintile (mean Q1 ceiling = 49.6\% of income).

The affordability cap binds for 3.8\% of indebted households ($n = 178$), but eliminates 22.8\% of total calculated debt service, indicating concentrated downward adjustments on high-debt outliers.
For 53 households, gross income falls below statutory minimum living expenses, leaving zero calculated repayment capacity.
Their debt balances are retained, as the data cannot separate legacy debt, informal borrowing, or non-compliant credit, resulting in these highly constrained households entering the model with zero active debt service flows.

Reconstructed debt servicing is evaluated against the \ac{FinScope} self-reported budget share allocation \cite{finscope2019}.
Among households holding at least one modeled credit product, debt service accounts for 6.07\% of monthly spending on clothing accounts and hire purchase, rising to 9.78\% when including mortgages, credit cards, and auto loans.
Standardizing the model to the same population and denominator yields a debt service share of 12.8\% of monthly outlay—exceeding the empirical benchmark, which includes broader product coverage.
This divergence reflects using statutory ceiling \acp{APR} instead of realized market rates.
Across indebted agents, the calculated debt-service-to-income ratio (\ac{DSTI}) exhibits a weighted mean of 10.9\% and a median of 3.3\% (the full-population mean is 4.8\%, reflecting the 56.6\% zero-debt baseline).

A known limitation in the expenditure split is accounting overlap: the \ac{NIDS} non-food aggregate used for discretionary expenditure includes reported clothing, hire purchase, and vehicle payments (median = 12.6\% of non-food outlays among debt-paying households).
Consequently, constructed servicing charges a flow already present in discretionary spending, overstating compressible outlays for indebted agents and slightly understating true relative debt service proportions.
Servicing-sensitive simulation outputs should therefore be interpreted as comparative upper bounds.

\subsection{Population Resampling}
\label{sec:resample}

The 10{,}841 weighted \ac{NIDS} households are resampled with replacement proportional to survey weights (seed: 20260818) to form a standard unweighted agent population of $N = 5{,}000$, sampled from 3{,}221 unique source records.
Downstream simulation statistics are calculated unweighted and reflect empirical distributional properties directly.

Figure~\ref{fig:pop-fidelity} illustrates structural alignment between the resampled agents and the weighted source data.
The per-capita income distribution matches closely across all spectrums (Kolmogorov-Smirnov distance = 0.016).
Categorical variables (income quintile, household size, province, primary income source) deviate by at most 1.4 percentage points across all 24 sub-categories.
Resampling slightly increases concentration: the per-capita income Gini coefficient rises by 0.015, moving from 0.651 in the weighted sample to 0.667 in the agent population.

\begin{figure}[htbp]
\centering
\includegraphics[width=\textwidth]{figures/fig_population_fidelity.pdf}
\caption{Distributional fidelity of the agent population. Each panel compares the unweighted 5{,}000 agents against the weighted NIDS source they are drawn from. Panels (a) and (b) cover the per-capita income distribution and its concentration; panels (c) to (f) cover the categorical attributes that gate behaviour and define the peer reference groups in the simulation.}
\label{fig:pop-fidelity}
\end{figure}

\subsection{Validation of the Synthetic Population}
\label{sec:data-validation}

We evaluate the synthetic population using eighteen diagnostic validation checks across five domains: internal consistency against \ac{NIDS} (Group A), matching fidelity to FinScope marginals (Group B), external benchmarks (Group C), debt service plausibility (Group D), and resample fidelity (Group E).
Table~\ref{tab:validation} summarizes the scorecard results: seventeen tests satisfy target tolerances, with only the external debt service level check exceeding tolerance bounds.

\begin{table}[H]
\centering
\caption{Data-layer validation scorecard for the constructed agent population. Group E compares the unweighted 5{,}000 synthetic agents against the weighted empirical source data.}
\label{tab:validation}
\begin{tabular}{L{4.9cm}L{3.6cm}rrc}
\toprule
\textbf{Check} & \textbf{Metric} & \textbf{Value} & \textbf{Tolerance} & \\
\midrule
A. Quintile shares $\approx 20\%$ & max $|$share $-0.20|$ & 0.004 & $\pm$2pp & \checkmark \\
A. Mean household size & weighted mean & 3.03 & 3.3 $\pm$0.5 \cite{statssa_cs_2016} & \checkmark \\
A. No impossible values & negative values & 0 & exact & \checkmark \\
\addlinespace
B. Banked (national) & vs FinScope & 0.4pp & $\pm$3pp & \checkmark \\
B. Store card & vs FinScope & 1.8pp & $\pm$3pp & \checkmark \\
B. Revolving credit & vs FinScope & 0.5pp & $\pm$3pp & \checkmark \\
B. Hire purchase & vs FinScope & 0.5pp & $\pm$3pp & \checkmark \\
B. Short-term loan & vs FinScope & 0.6pp & $\pm$3pp & \checkmark \\
B. Personal loan & vs FinScope & 0.2pp & $\pm$3pp & \checkmark \\
B. By-quintile flags & max over q$\times$flag & 2.3pp & $\pm$5pp & \checkmark \\
\addlinespace
C. Per-capita income Gini & weighted Gini & 0.651 & in (0.63, 0.70) & \checkmark \\
C. Aggregate debt vs \ac{CCMR} & stock $/$ in-scope book & 0.49 & in (0.30, 0.80) & \checkmark \\
\addlinespace
D. Servicing plausible pre-guard & debtors the guard binds & 3.8\% & $\leq$10\% & \checkmark \\
D. Debt service vs \ac{FinScope} & monthly outlay, debtors & 12.8\% & in (6.1, 9.8)\% & $\times$ \\
\addlinespace
E. Resample quintile shares & max $|$unwtd $-$ wtd$|$ & 1.3pp & $\pm$2pp & \checkmark \\
E. Resample flag rates & max $|$unwtd $-$ wtd$|$ & 0.8pp & $\pm$3pp & \checkmark \\
E. Resample income ECDF & KS max gap & 0.016 & $\leq$0.05 & \checkmark \\
E. Resample Gini & $|$agents $-$ source$|$ & 0.015 & $\leq$0.02 & \checkmark \\
\bottomrule
\end{tabular}
\end{table}

Group C assesses performance against macroeconomic benchmarks.
The per-capita income Gini of 0.651 falls within the expected target interval (0.630-0.700), bounded below by the World Bank's consumption-based Gini of 0.630 \cite{worldbank_gini_zaf} (income-based inequality measures in South Africa systematically exceed consumption-based estimates).
Total synthetic population debt (ZAR 191.4 billion) accounts for 49\% of the corresponding \ac{CCMR} gross debtor book (ZAR 392.0 billion) \cite{ncr_ccmr_2017}.
This ratio is consistent with standard self-reporting undercount gaps in micro-surveys relative to administrative registers.
This result supports our non-response zero-fill choice in Section~\ref{sec:backbone}: imputing missing debt entries at holder means would generate aggregate debt equal to 111\% of the entire administrative credit book.

Calculated debt service exceeds self-reported survey averages by roughly 3 percentage points of monthly outlay, aligned with expected upward bias from capping interest rates at statutory maximums.
Consequently, model results sensitive to baseline debt service levels are interpreted as conservative upper bounds.

The remaining diagnostic checks verify population construction fidelity. Two baseline structural deviations remain: mean household size (3.03) passes tolerance but sits below the 3.3 estimate from the 2016 Community Survey \cite{statssa_cs_2016}.
Because weighted population totals match national census benchmarks, this generates roughly 10\% more total household units (18.67 million vs 16.9 million aggregate units).
Additionally, Group B checks validate internal match consistency relative to donor distributions rather than external indicator accuracy.

Finally, the synthesized population exhibits baseline structural financial vulnerability prior to simulation runs.
At initialization, 3.9\% of households cannot cover committed expenditures from income, 5.2\% cannot cover committed outlays plus scheduled debt service, and 44.3\% hold zero liquid savings.
Under the default rules detailed in Section~\ref{sec:submodels}, these households register financial distress at $t=0$ independent of experimental shocks or \ac{BNPL} availability.

\subsection{Agent Attribute Mapping}
\label{sec:agent-mapping}

Table~\ref{tab:agent-mapping} summarizes the complete state variables initialized for each agent. Variables map directly to microdata observations, structural imputations, or derived accounting rules; unused survey fields are excluded.

\begin{table}[H]
\centering
\caption{Data-to-agent structural mapping. All monetary fields are denominated in 2017 ZAR. Flow variables are scaled from the monthly survey basis to the 14-day tick by $12/26$.}
\label{tab:agent-mapping}
\begin{tabular}{L{4.8cm}L{4.4cm}L{4.2cm}}
\toprule
\textbf{Agent field} & \textbf{Description} & \textbf{Source} \\
\midrule
\texttt{income\_monthly} & Total monthly liquidity inflows & NIDS \\
\texttt{income\_wage\_tick} & Wage component; the base a shock removes & NIDS \\
\texttt{n\_earners} & Employed members on the roster & NIDS individual file \\
\texttt{income\_source} & Dominant component: wage, grant, other & NIDS, derived \\
\texttt{committed\_tick} & Incompressible food and rent costs & NIDS, derived \\
\texttt{discretionary\_tick} & Non-food expenditure & NIDS, derived \\
\texttt{liquid\_savings} & Proxy for immediate cash buffer & NIDS, winsorised \\
\texttt{d\_trad} & Outstanding traditional debt balance & NIDS \\
\texttt{apr\_annual}, \texttt{term\_months} & Rate and horizon pricing that balance & Product mix (Section~\ref{sec:servicing}) \\
\texttt{scheduled\_service\_tick} & Monthly debt servicing obligation & Constructed (Section~\ref{sec:servicing}) \\
\texttt{nca\_capacity\_monthly} & Regulation 23A affordable residual & Statutory (Section~\ref{sec:servicing}) \\
\texttt{banked} & Eligibility gate for \ac{BNPL} & FinScope, imputed \\
\texttt{income\_quintile} & Relative wealth bracket; donor match key & NIDS, weighted \\
\texttt{province} & Geographic tag; donor match key & NIDS \\
\texttt{reference\_group} & Network topology for peer adoption observation & Derived (quintile $\times$ province) \\
\texttt{household\_size} & Total household members & NIDS \\
\bottomrule
\end{tabular}
\end{table}

Two simulation dimensions are unobserved in the microdata: all agents initialize with zero \ac{BNPL} balances to establish an unconfounded baseline, and individual psychological repayment archetypes (Section~\ref{sec:submodels}) are assigned stochastically at initialization.